% Source: source/普通高考/2014/2014重庆文.pdf, page 1, question 5.
% pdfimages -list found no embedded image objects; the source flowchart is vector artwork.
% Grid evidence: tmp/review_flowchart_2014_chongqing_liberal/grid/page-1-grid100.png
% (100px coarse + 50px fine) and q5-block-grid50.png (50px + 25px fine).
% Comparison crop: tmp/review_flowchart_2014_chongqing_liberal/crop/q5_original_final.png (no-grid page render).
% Program topology: 开始 -> k=2,s=0 -> k<10; 是 -> s=s+k -> k=2k-1 -> return to
% the shared stem above the test; 否 -> 输出 k -> 结束. All connectors are solid
% directed edges; branch labels are placed beside the corresponding exits.

\begin{tikzpicture}[
    >=Stealth,
    x=1cm,
    y=0.90cm,
    line width=0.45pt,
    line cap=round,
    line join=round,
    font=\normalsize,
    flowtext/.style={anchor=center, align=center,
        execute at begin node={\setlength{\parindent}{0pt}}},
    flowbox/.style={draw, flowtext, minimum height=0.68cm,
        inner xsep=5pt, inner ysep=3pt},
    terminal/.style={flowbox, rounded corners=0.16cm, minimum width=1.35cm},
    process/.style={flowbox, minimum width=2.95cm},
    decision/.style={flowbox, diamond, aspect=3.0, minimum width=3.70cm,
        minimum height=1.04cm},
    io/.style={flowbox, trapezium, trapezium left angle=72,
        trapezium right angle=108, minimum height=0.78cm},
]
    \path[use as bounding box] (-2.05,0.75) rectangle (4.90,8.30);

    \node[terminal] (start) at (0,7.95) {开始};
    \node[process, minimum width=3.10cm] (init) at (0,6.65) {$k=2,\ s=0$};
    \node[decision] (test) at (0,4.45) {$k<10$};
    \node[io, minimum width=2.05cm] (output) at (0,2.55) {输出 $k$};
    \node[terminal] (finish) at (0,1.20) {结束};

    \node[process, minimum width=2.50cm] (sum) at (3.20,4.45) {$s=s+k$};
    \node[process, minimum width=2.95cm] (update) at (3.20,5.45) {$k=2k-1$};

    % Main sequence and terminating branch.
    \draw[->] (start.south) -- (init.north);
    \draw[->] (init.south) -- (test.north);
    \draw[->] (test.south) -- node[left=2pt] {否} (output.north);
    \draw[->] (output.south) -- (finish.north);

    % Loop branch: test -> sum -> update, then a left-pointing return arrow
    % into the shared stem above the test. The vertical join has no duplicate
    % arrowhead; the main stem carries the single directed arrow into test.
    \draw[->] (test.east) -- node[above=4pt] {是} (sum.west);
    \draw[->] (sum.north) -- (update.south);
    \draw[->] (update.north) -- (3.20,6.15) -- (0,6.15);
    \draw (0,6.15) -- (0,5.10);
\end{tikzpicture}
